\newglossaryentry{alometria}{
	name=Alometria,
	description={relação de escala entre duas grandezas de um sistema na forma de lei de potência, $y = c\,x^{\beta}$; é linear se $\beta = 1$, sublinear se $\beta < 1$ e superlinear se $\beta > 1$}
}

\newglossaryentry{bacia}{
	name=Bacia hidrográfica,
	description={região cujo escoamento superficial converge para uma mesma saída; neste trabalho, o conjunto de pixels de terra que drenam para um mesmo \textit{sink} costeiro}
}

\newglossaryentry{foz}{
	name=Foz,
	description={ponto em que uma bacia hidrográfica deságua no oceano; aqui, o pixel-\textit{sink} da bacia na grade de relevo}
}

\newglossaryentry{mancha}{
	name=Mancha urbana,
	description={aglomerado de pixels densamente povoados e próximos entre si, definido pelo CCA sem recorrer a fronteiras administrativas}
}

\newglossaryentry{nurdle}{
	name=\textit{Nurdle},
	sort=nurdle,
	description={\textit{pellet} de resina plástica pré-produção, com alguns milímetros de diâmetro, usado como matéria-prima da indústria de transformação de plásticos}
}

\newglossaryentry{percolacao}{
	name=Percolação invasiva,
	description={processo em que um fluido invade um meio sempre pelo ponto de menor resistência da fronteira já invadida; no relevo, a resistência é a altitude}
}

\newglossaryentry{sink}{
	name=\textit{Sink},
	sort=sink,
	description={pixel de saída de uma bacia na percolação; todo pixel de terra vizinho ao oceano é um \textit{sink}}
}

\newglossaryentry{zipf}{
	name=Lei de Zipf,
	description={regularidade empírica segundo a qual a população da $r$-ésima maior cidade decai como $r^{-1}$}
}
